% =====================================================================
\chapter{Conclusion and Future Work}
\label{ch:conclusion}
% =====================================================================

\section{Summary}

This dissertation studied whether hallucinations in large language models can be detected
cheaply by reading the model's own internal representations, and which of the many proposed
internal signals are worth their cost. We adopted the label-free, per-model data generation
of the MIND framework, kept its canonical last-token, last-layer hidden state as a backbone,
and added three inexpensive signals from the same forward pass: how much the representation
drifts between layers, how much it varies across layers, and how uncertain the model is while
generating. We placed these in a controlled ablation alongside literature features and more
speculative features, fixed a single classifier and a single evaluation protocol, and
compared everything against four established baselines and two simple reference methods on a
suite of up to ten benchmarks --- all on free cloud hardware.

The honest summary of the findings is threefold. First, detection of held-out
in-distribution text works at a useful level, and better at larger data scales, but
transferring an encyclopedia-trained probe to question answering, summarisation, and dialogue
is hard: on transfer, every method we tested sits close to chance once a single unusually
separable benchmark is set aside. Second, on the one model for which we have a complete set
of file-backed measurements, the three core features add little on their own, the gains come
from richer feature combinations, and even the best configurations match rather than beat the
strongest baseline --- a supervised probe whose advantage seems to come from reading a better
layer rather than from a cleverer feature. Third, the result is model-dependent: on a second
model of comparable size a proposed configuration --- built from the more exploratory
internal-state signals --- is instead the strongest method of all, clearing every baseline.
Two models of similar size disagreeing about which method wins is itself the finding, and a
caution against reading a universal winner off any one model.

The contribution is therefore not a new state-of-the-art detector. It is a unified,
free-tier-reproducible comparison of internal-state hallucination features across several open
models, built so that the regime-dependent and model-dependent effects above are visible and
checkable rather than hidden inside a single headline number.

\section{Limitations}

We state the limitations plainly, because several of them bound how far the conclusions can
be taken.

\textbf{Two models, not the whole fleet.} The comparison rests on two six-billion-parameter
models: one with the complete sixteen-configuration, six-baseline sweep, and one with the six
baselines measured against the exploratory-feature configurations. The remaining large models
have their training data ready but their detection runs are still in progress, so the
cross-model claim is drawn from two models rather than the full fleet.

\textbf{Modest sample sizes.} The unified runs use two thousand training examples per model
and evaluate on a few hundred rows per benchmark. At this scale, individual AUROCs are noisy,
classifiers sometimes collapse to a single class, and small differences between methods should
not be over-interpreted. The earlier, larger-scale in-distribution figures suggest that more
data helps, but we did not re-run the full ablation at that scale.

\textbf{Free-tier constraints.} The nine-hour session limit, the weekly GPU budget, and the
sixteen-gigabyte cards shaped the design: some baselines were dropped for cost, evaluation was
capped at a fixed number of rows, and one gated model could not be run at full data size.
These are reasonable trade-offs for a reproducible free-hardware study, but they are
constraints, and a better-resourced replication might reach steadier numbers.

\textbf{One benchmark dominates.} The transfer averages on the fully measured model are
carried by a single easy benchmark. The rankings should be read with that in mind, and
per-benchmark results (Appendix~\ref{app:tables}) are more informative than the averages.

\textbf{Split evaluation on the second model.} On the second model the literature and
core-feature configurations were scored on the earlier seven-benchmark suite while the
exploratory-feature configurations and the baselines used the full ten; a single unified
re-run would put all sixteen configurations on the same footing.

\section{Future work}

Several concrete next steps follow from these limitations.

The most immediate is to \emph{complete the sweep}: run the full sixteen-configuration,
six-baseline comparison on the second model and on the remaining seven-billion models, and
regenerate the gated model's data at full size. This broadens a two-model comparison into a
fleet-wide one and would settle whether the model-dependence seen here is the rule.

The discussion pointed at \emph{layer selection} as the lever most worth pulling. The strongest
baseline's edge came from reading a mid-network layer rather than the final one, so a focused
ablation over the layer at which the canonical and internal-state features are read --- on each
model --- is likely to be more productive than adding further scalar features.

A third direction is \emph{scale and calibration}. Re-running the ablation at the larger
data scale used in the mid-term evaluation would test whether the small per-feature effects
become reliable, and reporting calibration error and Brier score alongside AUROC would make the
detectors more useful as decision-support tools, since a calibrated probability is worth more
than a bare verdict.

Finally, the finding that transfer is the hard part suggests studying the \emph{training
distribution} itself: mixing a small amount of task-like data into the encyclopedia
continuations, or generating continuations from more varied sources, may close part of the
in-distribution-to-transfer gap that dominates the present results.

\section{Closing remark}

The modest headline of this work --- that inexpensive internal-state features are competitive
but not decisive, and that cross-task transfer remains hard --- is, we think, more useful than a
louder claim would have been. The detectors in this area are often reported under favourable,
inconsistent conditions; holding everything fixed and running them on free hardware shows where
the real difficulty lies. Making that difficulty visible, and reproducible, is the point.
